\subsection{What rises: the parents' wage, their wealth and the son's skill}

The parents' wage and wealth carry most of what is measured in the IGE, and the son's skill at entry carries most of
its rise with age. Writing the IGE as the coefficient $\beta$ of the son's log earnings on log parental income, and the
same regression with parental components $m_k$ added as $\tilde\beta$ and $\gamma_k$, the split
$\beta = \tilde\beta + \sum_k \gamma_k \delta_k$ assigns to each component the part of the IGE carried by
parental-income differences that are differences in it, where $\delta_k$ is the component's gradient in parental income
\citep{gelbach2016covariates}. Table~\ref{tab:split_main} reports it on all sons and on the sons born 1964 and later,
whose teens fall in the PSID's wealth waves. What holds in every version of the split: the parents' hourly wage carries
a fifth or more of the IGE (20.5\%$^{***}$ on the later cohorts; 37 percent once measurement error in the wage is
corrected); the parents' net worth, where it is measured, carries between a seventh and three tenths, depending on how
it enters (14.5\%$^{***}$ as the inverse hyperbolic sine, 29.1\%$^{***}$ as a rank), its contribution doubling from
0.062$^{**}$ at 25--29 to 0.127$^{***}$ at 35--39, and most of it is housing; the parents' schooling carries an eighth without wealth
and less with it; spells out of work carry a few percent; hours and the state's labor market nothing. With wealth in,
little of the IGE is left as parental income at equal measures (3.5 percent); on all sons, without wealth, half
is.\footnote{The full split---the original components and the rest, the four-by-four table of what passes to what, the
same split in four data sets, the measurement-error correction and what identifies the rest---is in
Appendix~\ref{sec:passedon}. Code: \texttt{75.WhatIsPassedOn.py}, \texttt{80.WhatIsPassedOnRest.py},
\texttt{96.SplitMeasurementError.py}, \texttt{97.RestIdentification.py}, \texttt{98.ParentalWealth.py}.} The split is
accounting: the components are correlated choices and outcomes.

\begin{table}[!htbp]
\centering
\caption{The components of the IGE: shares of the IGE of sons' earnings (PSID)}
\label{tab:split_main}
\resizebox{\ifdim\width>\textwidth\textwidth\else\width\fi}{!}{%
\resizebox*{!}{\ifdim\totalheight>0.85\textheight 0.85\textheight\else\totalheight\fi}{%
\begin{tabular}{lcccc}
\toprule
 & \begin{tabular}[b]{@{}c@{}}All sons, 25--55\end{tabular} & \begin{tabular}[b]{@{}c@{}}Born 1964+, 25--29\end{tabular} & \begin{tabular}[b]{@{}c@{}}Born 1964+, 35--39\end{tabular} & \begin{tabular}[b]{@{}c@{}}Born 1964+, 25--55\end{tabular} \\
 & (1) & (2) & (3) & (4) \\
\midrule
\multicolumn{5}{l}{\textit{Share of the IGE (\%)}} \\
Parents' hourly wage & $12.0^{**}$ & $22.0^{*}$ & $20.6^{**}$ & $20.5^{***}$ \\
 & (5.9) & (11.5) & (8.3) & (6.4) \\
Parents' schooling & $12.9^{***}$ & $-5.4$ & $9.4^{*}$ & $3.9$ \\
 & (3.7) & (6.4) & (5.5) & (4.5) \\
Parents' net worth & --- & $23.0^{***}$ & $31.2^{***}$ & $29.1^{***}$ \\
 &  & (8.7) & (8.0) & (7.0) \\
Parents' spells out of work & $4.9^{**}$ & $8.9^{*}$ & $4.9$ & $6.8^{**}$ \\
 & (2.2) & (4.6) & (4.5) & (3.4) \\
Parents' hours & $1.2$ & $3.8$ & $4.4$ & $2.9$ \\
 & (1.4) & (3.2) & (2.7) & (2.3) \\
State labor market & $0.2$ & $1.4$ & $0.9$ & $1.7$ \\
 & (0.5) & (1.2) & (1.3) & (1.2) \\
Transfers (AFDC/TANF years, transfer share) & $6.5^{*}$ & $5.2$ & $9.7$ & $7.9$ \\
 & (3.8) & (7.2) & (7.8) & (5.7) \\
Asset income, home ownership & $0.1$ & $2.3$ & $-6.2$ & $-4.5$ \\
 & (2.7) & (6.4) & (5.2) & (4.6) \\
Labor income beyond wage and hours & $8.9$ & $-1.6$ & $10.9$ & $12.1$ \\
 & (7.5) & (25.4) & (19.9) & (12.8) \\
Other members' income & $0.3$ & $2.8$ & $4.7$ & $4.1^{**}$ \\
 & (1.1) & (2.7) & (2.9) & (1.9) \\
Family size & $0.1$ & $0.3$ & $-0.8$ & $-0.2$ \\
 & (0.7) & (1.1) & (0.9) & (0.7) \\
Rest: parental income at equal measures & $47.6^{***}$ & $15.5$ & $3.1$ & $3.5$ \\
 & (10.8) & (33.8) & (24.2) & (16.4) \\
IGE of ln earnings & $0.396^{***}$ & $0.271^{***}$ & $0.408^{***}$ & $0.384^{***}$ \\
 & (0.024) & (0.025) & (0.042) & (0.031) \\
\bottomrule
\end{tabular}}}
\par\smallskip
\begin{minipage}{\textwidth}\footnotesize \textit{Notes:} Exact Gelbach split of the IGE of sons' ln earnings into parental components (means over the son's ages 13--19) and the rest, parental income at equal measures. Column 1: all sons, without wealth (the PSID measures net worth only from 1984). Columns 2--4: sons born 1964 and later, whose teens fall in the wealth waves, with parental net worth (rank within the son's birth year, mean over the wealth waves at 12--20). What holds in every version (Appendix~B: the original split, the measurement-error correction, four data sets): the parents' wage carries a fifth or more, wealth three tenths where it is measured, the parents' schooling an eighth without wealth and less with it, spells a few percent, hours and place nothing. PSID weights; 200 cluster-bootstrap draws of 1968 families. Source: \texttt{80.WhatIsPassedOnRest.py}, \texttt{98.ParentalWealth.py}; table built by \texttt{100.PaperBFableTables.py}. $^{***}$ $p<0.01$, $^{**}$ $p<0.05$, $^{*}$ $p<0.1$.\end{minipage}
\end{table}


The rise of the IGE with age is mostly the son's skill at entry paying more later. In the NLSY79 the earnings IGE rises
from 0.340$^{***}$ at 25--29 to 0.506$^{***}$ at 50--54, and the son's AFQT, measured before he enters the labor market,
carries 0.103$^{***}$ of that rise, three fifths; the parents' schooling beyond it carries most of the rest. It is age,
not the rise of the skill premium: in the NLSY97, whose sons reach the same ages twenty years later, the AFQT's
contribution rises by the same factor (1.64 against 1.70 from 25--29 to 40--44), although it carries a quarter of the
NLSY97 opening rather than three fifths (panel B of Table~\ref{tab:age_period}). Because the returns to both schooling
and measured skill rise with age, the pattern is that of skill complementing the human capital accumulated on the job
\citep{bhuller2017life}, not of employers learning about it \citep{altonji2001employer, lange2007speed}; in the PSID the
rise is within occupations.\footnote{Appendix~\ref{sec:skill} has the AFQT split, the occupation decomposition of the
rise, the early job and the return to schooling. Code: \texttt{86.NLSY79YouthWhatIsPassedOn.py},
\texttt{87.NLSY79YouthAFQT.py}, \texttt{94.NLSY97AFQT.py}.}

One component's share rises with age: the parents' spells out of work, from 0.017$^{**}$ of the IGE at 25--29 to
0.060$^{***}$ at 45--49, 13.5\%$^{***}$ of it. It is the father's layoffs, and Sections~\ref{sec:stability}--\ref{sec:daughters}
are about them. The son's own spells run the other way: they carry 30.9\%$^{***}$ of the earnings IGE at 25--29 and
8.7\%$^{***}$ at 45--49, so that the IGE of a young man is in large part about whether he works and that of a man in
mid-career about what he is paid; net of them, the life-cycle bias of an IGE measured at 25--29 is larger (Appendix~\ref{app:spells}).

